\section{Discussion}\label{sec:discussion}

\subsection{Principal findings}
This study shows that the routine clinicopathological record carries moderate, reproducible
information about breast cancer grade. Macro-AUCs were \IntAucMinR--\IntAucMaxR{} internally and
\ExtAucMin--\ExtAucMax{} across five external cohorts from three continents. Grade~3 was the easiest
class to recognise and grade~2 the hardest, which mirrors the known difficulty of grade~2 for
pathologists~\citep{boiesen2000grading,vandooijeweert2020variation}. Hormone-receptor status was the
most informative variable, and tumour size, nodal status, HER2 and lobular histology contributed in
the directions that clinicians would expect.

Three further findings matter for anyone who builds or uses such models. First, flexibility did not
help: no flexible learner beat multinomial logistic regression internally, and external differences
were rare and went in both directions. This agrees with systematic reviews of clinical prediction
models~\citep{christodoulou2019lr} and is unsurprising when the predictors are few and mostly binary.
Second, discrimination and predictor effects transported well (calibration slopes close to one), but
the average level of the predictions did not. The same frozen model under-predicted grade~3 in
MSK-2018 and over-predicted it in Duke, because the share of grade~3 tumours differs more between
cohorts than their predictors do. A simple intercept update fixed the problem in both cohorts.
Third, the size of the leakage effect was striking. Each shortcut that is easy to make with public
breast cancer data inflated performance, by up to \LeakMaxInfl{} in macro-AUC. That is more than the
whole distance between a useless model and the best correctly validated one.

\subsection{What the imaging extension adds}
In the Duke cohort, 529 MRI radiomic features did not improve on the routine record, and adding them
lowered cross-validated discrimination (paired $\Delta$ \DukeFusDelta, \DukeFusDeltaLo{} to
\DukeFusDeltaHi). With \NDukeRad{} graded patients and hundreds of correlated features, the model
could not find signal beyond what receptor status, tumour size and nodal status already provide.
Published radiomics results for grade~\citep{wang2022duke,wang2026dcemrigrade} should therefore be
judged against a clinical baseline evaluated in the same patients, as done here.

\subsection{What this means for Saudi Arabia}
Re-weighting to the Saudi profile raised the grade-3 share of every Western cohort by \ShiftMin{} to
\ShiftMax{} percentage points. The Saudi case-mix---fewer receptor-positive and more HER2-positive
tumours---is therefore associated with more high-grade disease. The level reached, however,
depended on the source cohort (\RwGThreeKMin--\RwGThreeKMax\% for the Dammam profile), because the
Western cohorts already differed in their grade-3 share (\ObsGThreeMin--\ObsGThreeMax\%). Published
Saudi series report less grade~3 than the METABRIC- and MSK-based projections: \SaudiGThreeA\% in
Dammam, \SaudiGThreeB\% in Riyadh and \PoolGThree\% when pooled across \PoolK{} series. A model
developed in METABRIC would predict, on average, \PredGThreeMet\% grade~3 for patients with the
Dammam profile, about \OverPredPP{} percentage points too many. Several explanations are possible:
differences in grading practice (grade~2 is the least reproducible category), referral and era
effects in hospital series, or real differences in tumour biology. Summary data cannot separate them.
The practical message is clear, though. A Western grade model should not be used for Saudi patients
without recalibrating its intercept, which needs only the local grade distribution. The most valuable
next step is an openly shared, de-identified, grade-labelled Saudi cohort; the code released with this
paper can validate and recalibrate the models on such a cohort with one additional data-mapping
function.

\subsection{Strengths and limitations}
The main strengths are full reproducibility from public data with recorded checksums, an analysis plan
fixed before the experiments, five external cohorts from three continents (one recruited through
imaging rather than genomics), paired model comparisons, and an explicit measurement of leakage.

The study also has important limitations. Grade was assigned by different pathologists, in different
eras, under slightly different conventions; the Ljubljana ``degree of malignancy'' predates the
Elston--Ellis modification. Disagreement between pathologists~\citep{boiesen2000grading,longacre2006interobserver}
limits the discrimination that any model can reach. Several cohorts were assembled for genomic
studies or enriched for advanced disease, so they are not population-based. Tumour size was reduced
to T category, and tumour size and nodal status were missing in three external cohorts. The
sample-size criteria were \SSMet: the grade~1 versus grade~2 comparison, which carries little signal,
would need \SSRequired{} patients against \SSAvailable{} available. Penalised estimation and
calibration slopes close to one suggest that overfitting was small in practice, but estimates at the
grade~1/2 boundary should be read with care. The Saudi analysis relies on published summary figures,
not patient data. The Dammam series does not report the share of patients aged $\le$50 years (the
results were not sensitive to the assumed value); the Riyadh series comes from a referral centre with
many stage~IV patients and about 10\% non-Saudi patients; the Madinah series grades ductal
carcinoma only; and the Najran series grades all lesions, including in-situ and benign ones, so it
was used only in a sensitivity analysis. \TabpfnLimitation{}Finally, the image-level analyses that
we had planned with the Saudi KAU-BCMD mammograms~\citep{alsolami2021kau} and the Egyptian CDD-CESM
and BUSI collections~\citep{khaled2022cddcesm,aldhabyani2020busi} could not be run, because their
hosts were not reachable from the build environment.

\subsection{Implications for practice and research}
A routine-variable model is not a substitute for histological grading. It has three realistic uses:
imputing missing grade in registry or trial data with calibrated uncertainty; flagging records whose
assigned grade is unlikely given the rest of the record, as a quality-assurance check; and serving as
the minimum baseline that image-based grading tools must beat in the same patients, with
patient-level data splits and with grade-encoding columns removed. Each use requires local
recalibration of the intercept.

\section{Conclusions}
This study answers a simple question with open, reproducible evidence: routine clinicopathological
variables predict breast cancer grade moderately and consistently, with macro-AUCs of about
\IntAucMaxR{} internally and up to \ExtAucMax{} in external cohorts from three continents.

The results contribute to knowledge in four ways. First, they set a transparent, externally
validated benchmark for grade prediction that image-based methods can be measured against. Second,
they show that logistic regression is enough for this task; gradient boosting, random forests and MRI
radiomics added nothing. Third, they show that predictor effects transport across countries but
baseline grade frequencies do not, so calibration, not discrimination, is the weak point when a
grade model moves to a new hospital. Fourth, to our knowledge they are the first to quantify how much common
leakage shortcuts inflate apparent performance in grade prediction---by up to \LeakMaxInfl{} in
macro-AUC---which gives reviewers and readers a concrete reason to distrust near-perfect results.

For practitioners, the message is practical. A simple, interpretable regression model can impute
missing grade or flag implausible grades, provided it is recalibrated to the local grade
distribution. Models that use the Nottingham Prognostic Index or the component scores, or that
preprocess data before splitting, should not be trusted. For Saudi Arabia, the results show that the
local case-mix implies more high-grade disease than in Western cohorts, yet Western-trained models
would still over-predict grade~3 compared with published Saudi series. Local recalibration is
therefore essential, and an openly shared, grade-labelled Saudi cohort would be the most valuable
next contribution. All code, settings and results are public, so these findings can be checked,
updated and extended as new data become available.
